\documentclass{article}

\IfFileExists{neurips_2025.sty}{
  \usepackage[preprint]{neurips_2025}
}{
  \usepackage[margin=0.78in]{geometry}
}

\usepackage[utf8]{inputenc}
\usepackage[T1]{fontenc}
\usepackage{microtype}
\usepackage{booktabs}
\usepackage{array}
\usepackage{amsmath}
\usepackage{graphicx}
\usepackage{hyperref}
\usepackage{enumitem}
\usepackage{xcolor}

\hypersetup{
  colorlinks=true,
  linkcolor=blue,
  citecolor=blue,
  urlcolor=blue
}

\setlist[itemize]{leftmargin=*,topsep=2pt,itemsep=1pt}
\setlist[enumerate]{leftmargin=*,topsep=2pt,itemsep=1pt}
\newcommand{\metric}[1]{\textbf{#1}}

\title{Revenue Follows Seasonal Demand; COGS Reflects Margin and Promotion Regime}

\author{
Datathon 2026 Round 1 Team \\
\texttt{<GitHub repository link>}
}

\begin{document}
\maketitle

\begin{abstract}
We forecast daily Revenue and COGS for an e-commerce business from 2023-01-01 to 2024-07-01 using only the provided historical datasets and calendar structure. The central business finding is that Revenue is primarily governed by demand seasonality and regime-level conversion, while COGS is better explained through a margin-ratio mechanism driven by month, discount pressure, and historical H2 instability. Our final forecasting pipeline combines a stable demand anchor with a separate COGS ratio head and a conservative train-derived daily allocation layer. The final forecast achieved a Kaggle MAE of \metric{667,139.05}.
\end{abstract}

\section{Problem Context and Data Provenance}

The task is to forecast one row per day with columns \texttt{Date}, \texttt{Revenue}, and \texttt{COGS} for the horizon 2023-01-01 to 2024-07-01. The evaluation emphasizes low \metric{MAE}, low \metric{RMSE}, and high \metric{$R^2$}, while the report rubric also rewards clear EDA, time-aware validation, feature explanation, leakage control, and reproducibility.

\paragraph{Business framing.}
The forecast is not only a statistical exercise. Revenue determines demand planning and cash-flow expectations, while COGS determines margin risk and inventory-promotion tradeoffs. We therefore model Revenue and COGS with different business mechanisms: demand level and seasonality for Revenue, margin and promotion regime for COGS.

\paragraph{Target reconstruction.}
The historical target identity is stable: Revenue is reconstructed as \texttt{quantity * unit\_price}, and COGS is reconstructed from \texttt{quantity * products.cogs}. The reconstruction audit over 3,833 training days had near-zero Revenue error and COGS error only at rounding scale.

\begin{table}[h]
\centering
\caption{Target reconstruction audit on training data.}
\label{tab:target_identity}
\begin{tabular}{lrr}
\toprule
Target & Max absolute error & Mean absolute error \\
\midrule
Revenue & $1.86 \times 10^{-9}$ & $2.33 \times 10^{-11}$ \\
COGS & $0.0050$ & $0.0025$ \\
\bottomrule
\end{tabular}
\end{table}

\paragraph{Feature availability.}
All modeling features and priors are derived from the provided historical datasets and calendar structure. The forecasting pipeline is reproducible from the submitted source code and does not use hidden target values or external datasets as model inputs. Operational source tables end at or before 2022-12-31, so future operational variables are only used through historical priors, not as future realized values.

\begin{table}[h]
\centering
\caption{Selected source coverage used for leakage-safe feature policy.}
\label{tab:source_coverage}
\begin{tabular}{lrrl}
\toprule
File & Rows & Date coverage & Role in model \\
\midrule
\texttt{sales.csv} & 3,833 & 2012-07-04 to 2022-12-31 & target history \\
\texttt{orders.csv} & 646,945 & 2012-07-04 to 2022-12-31 & funnel history \\
\texttt{web\_traffic.csv} & 3,652 & 2013-01-01 to 2022-12-31 & traffic priors \\
\texttt{promotions.csv} & 50 & 2013-01-31 to 2022-12-31 & recurring promo priors \\
\texttt{inventory.csv} & 60,247 & 2012-07-31 to 2022-12-31 & stockout priors \\
\bottomrule
\end{tabular}
\end{table}

\section{Business EDA: Demand Regime and Margin Pressure}

\paragraph{Insight 1: the structural break is a funnel break.}
The most important historical shift begins around 2019. Recent 2019-2022 Revenue is much lower than pre-2019 high-regime years, but sessions increased. The loss is therefore not a simple traffic problem: conversion and order volume collapsed, while AOV rose. This suggests a demand-quality regime shift.

\begin{table}[h]
\centering
\caption{Recent-low regime versus pre-2019 high regime. Ratios below 1 indicate decline.}
\label{tab:funnel_bridge}
\begin{tabular}{lrrrrr}
\toprule
Half & Revenue & Orders & Sessions & Conversion & AOV \\
\midrule
H1 & 0.596 & 0.484 & 1.318 & 0.365 & 1.235 \\
H2 & 0.570 & 0.456 & 1.301 & 0.345 & 1.248 \\
\bottomrule
\end{tabular}
\end{table}

\paragraph{Business interpretation.}
Traffic growth without order conversion implies that acquisition volume alone is not enough. The forecast must account for lower buying intent and a higher-ticket basket rather than extrapolating from sessions or calendar patterns only.

\paragraph{Insight 2: H1 is more seasonal; H2 is more unstable.}
Historical daily-shape correlations show that H1 seasonality is more repeatable than H2. Pairwise daily-shape correlation is 0.815 in H1 and 0.708 in H2. The model therefore uses stronger shrinkage and more conservative allocation in H2, especially around August and promotion-heavy months.

Exact lunar-calendar diagnostics show that Tet is a measurable but secondary seasonal effect. The first days of Tet are slightly above the local monthly baseline, while the broader pre/post-Tet window is noisier. This makes lunar calendar useful as a derived calendar feature, but not strong enough to replace the main demand-regime and margin-ratio logic.

\paragraph{Insight 3: COGS is a margin-ratio problem.}
COGS should not be treated as a passive follower of Revenue. The COGS/Revenue ratio has much higher dispersion in H2, and August is the most volatile month. This supports a separate ratio head:
\[
  \text{COGS}_t = \text{Revenue}_t \times \widehat{\text{COGS ratio}}_t.
\]

\begin{table}[h]
\centering
\caption{COGS/Revenue ratio dispersion from training history.}
\label{tab:ratio_dispersion}
\begin{tabular}{lrrrr}
\toprule
Segment & Mean & Std. dev. & P10 & P90 \\
\midrule
H1 & 0.826 & 0.050 & 0.782 & 0.909 \\
H2 & 0.930 & 0.159 & 0.788 & 1.050 \\
August & 1.083 & 0.287 & 0.785 & 1.421 \\
\bottomrule
\end{tabular}
\end{table}

\paragraph{Insight 4: promotions explain margin pressure.}
Promotion and discount variables dominate the COGS-ratio correlation table. The strongest Spearman correlations with daily COGS ratio are discount value mean (0.830), promo line share (0.814), average discount rate (0.810), promo duration (0.804), and active promo count (0.798). Since future promotions are not directly known, the model uses recurring train-derived promo priors by month and month-day.

\paragraph{Secondary mix findings.}
Product and geography mix shifts support the same story. Streetwear and Balanced segments gained share, Outdoor declined, Central region gained share, and West declined. These are useful for business explanation and sanity checks, while the final daily forecast remains governed by demand level, calendar allocation, and margin ratio.

\paragraph{Operational friction findings.}
The converted analyst document reinforces several business-facing signals that are consistent with the clean EDA logs. Organic search appears large in volume but weaker in revenue per session, pointing to a traffic-quality issue rather than a traffic-volume issue. Streetwear is both the dominant revenue category and the largest stockout-risk area. Returns are driven more by fit and expectation mismatch than delivery speed, with wrong-size returns being the most visible leakage driver. COD also creates order-quality risk because its cancellation rate is substantially higher than prepaid methods.

\section{Forecasting Method}

The modeling design follows the EDA: Revenue and COGS are related but not identical business processes. Revenue is a demand process; COGS is a margin process.

\paragraph{Revenue model.}
Revenue is decomposed into a period-level demand anchor and a daily seasonal allocation:
\[
  \widehat{\text{Revenue}}_t =
  \widehat{\text{Demand level}}_{\text{period}(t)}
  \times
  \widehat{\text{Seasonal share}}_t.
\]
The demand anchor uses historical Revenue, calendar recency, and train-derived regime priors. The daily allocation uses month, day-of-week, month-day, boundary-day, and historical seasonal profiles. H1 receives more shape confidence because its historical pattern is more stable; H2 is regularized more strongly.

\paragraph{COGS ratio head.}
Instead of directly scaling COGS from Revenue, the model estimates a COGS/Revenue ratio prior from historical month, discount, promotion, and regime evidence:
\[
  \widehat{\text{COGS}}_t =
  \widehat{\text{Revenue}}_t
  \times
  \widehat{r}_t,\quad
  r_t = \text{COGS}_t / \text{Revenue}_t.
\]
This separates margin pressure from demand level and makes the model robust to the fact that promotion-heavy days can have normal or high Revenue but unusually compressed gross margin.

\paragraph{Daily allocation and stabilization.}
The final configuration uses a conservative month-preserving daily stabilizer for H2 COGS ratio risk. It adjusts daily COGS allocation toward recent even-year historical ratio patterns while preserving period/month totals. This is intended to reduce tail-day error and avoid overreacting to unstable H2 daily spikes.

\paragraph{Feature policy.}
Features are grouped by availability:
\begin{itemize}
  \item \textbf{Known future:} calendar, month, day-of-week, month boundaries, and algorithmic lunar-calendar positions derived from \texttt{Date}.
  \item \textbf{Historical priors:} recurring promotion intensity, historical ratio profiles, demand-regime priors, source/category/region mix priors.
  \item \textbf{Excluded as realized future inputs:} future orders, realized traffic, future returns, future inventory states, and any unavailable post-2022 operational measurements.
\end{itemize}

\paragraph{Explainability.}
The key model drivers are interpretable by construction: demand regime, calendar seasonality, recurring promotion pressure, and COGS/Revenue ratio. For tree-based anchor components, feature importance and partial dependence checks are used as equivalent explanations to verify that the model relies on calendar, recency, and train-derived promotion/ratio signals rather than unavailable future information.

\section{Business Implications, Forecast Output, and Limitations}

\paragraph{From forecast to business action.}
The model is useful because it translates historical patterns into forward-looking business risks. The main forecast result is not only a Kaggle score; it is a structured view of where the business should pay attention: conversion quality, seasonal demand planning, H2 margin volatility, Streetwear inventory exposure, wrong-size returns, and COD cancellation risk.

\begin{table}[h]
\centering
\caption{Business findings translated into recommended actions and monitoring metrics.}
\label{tab:business_actions}
\begin{tabular}{p{0.25\linewidth}p{0.32\linewidth}p{0.28\linewidth}}
\toprule
Finding & Recommended action & Metrics to monitor \\
\midrule
Conversion declined despite higher sessions & Audit landing pages, product pages, cart flow, and checkout completion before increasing traffic spend & Conversion proxy, checkout completion, revenue per session \\
Organic search has high volume but weak commercial efficiency & Improve keyword-intent matching and landing-page relevance; retarget high-intent browsers & Organic revenue per session, organic conversion proxy \\
Revenue peaks in April--June & Prepare inventory, marketing, and logistics capacity before seasonal peaks & Monthly demand forecast, stockout days, fulfillment SLA \\
Streetwear dominates revenue and stockout exposure & Prioritize replenishment for high-demand Streetwear SKUs, sizes, and colors & Fill rate, days of supply, sell-through rate \\
COGS ratio is volatile in H2 and promotion-heavy months & Plan margin buffers and evaluate promotions by gross margin, not only revenue uplift & COGS/Revenue ratio, promo margin, discount rate \\
Wrong-size returns create refund leakage & Improve size guides, fit recommendations, and product descriptions & Wrong-size return count, refund amount, return rate \\
COD has higher cancellation risk & Confirm COD orders, limit COD on high-risk baskets, and incentivize prepaid payment & COD cancellation rate, prepaid share, cancelled value \\
\bottomrule
\end{tabular}
\end{table}

\paragraph{Forecasting output.}
The final forecast achieved a Kaggle MAE of \metric{667,139.05}. More importantly, the forecast is built around two business equations:
\[
  \widehat{\text{Revenue}} =
  \widehat{\text{Demand level}}
  \times
  \widehat{\text{Seasonal allocation}},
  \qquad
  \widehat{\text{COGS}} =
  \widehat{\text{Revenue}}
  \times
  \widehat{\text{Margin ratio}}.
\]
This framing is important because a revenue-only forecast can hide margin compression. A promotion-heavy day may look healthy in Revenue while still carrying high COGS risk.

\paragraph{Validation and sanity checks.}
Validation is time-aware. The workflow uses long-horizon and year-based folds to mimic forecasting more than one year ahead, rather than random folds that would leak adjacent temporal structure. Candidate forecasts are checked for row count, date range, non-negative values, period totals, maximum daily values, and COGS/Revenue ratio ranges before submission.

\paragraph{Limitations.}
Future operational facts such as actual 2023--2024 orders, returns, inventory, reviews, and traffic are unavailable at forecast time. The model therefore uses historical priors for these drivers. This makes the pipeline reproducible and avoids unavailable future inputs, but it limits the ability to react to new campaigns, supply shocks, or acquisition changes after 2022.

\paragraph{Summary.}
The recommended strategy is not to increase traffic or discounts indiscriminately. The business should first convert existing traffic more effectively, protect the core Streetwear category before seasonal peaks, and manage margin risk during promotion-heavy H2 periods. The forecast supports this strategy by separating demand seasonality from COGS ratio behavior.


\clearpage
\appendix
\section{Additional Evidence and Reproducibility}

\subsection{Requirement Checklist}
\begin{itemize}
  \item Forecast horizon: 2023-01-01 to 2024-07-01.
  \item Submission columns: \texttt{Date}, \texttt{Revenue}, \texttt{COGS}.
  \item Main metrics: MAE, RMSE, and $R^2$.
  \item Report format: NeurIPS-style LaTeX, 4 pages excluding references and appendix.
  \item Required methodology: time-aware validation, leakage control, explainability, and reproducible source code.
\end{itemize}

\subsection{Recurring Promotion Priors}
Training data shows recurring promotion activity by month. December and September contain full-month promotion coverage in historical data; July averages 23 active promotion days; August is high intensity but unstable. These are used as historical priors, not as future realized promotion schedules.

\begin{table}[h]
\centering
\caption{Selected recurring promotion month priors from training data.}
\begin{tabular}{lrrr}
\toprule
Month & Years with promo & Active days mean & Avg. discount max \\
\midrule
12 & 10 & 31.0 & 20.0 \\
9 & 10 & 30.0 & 11.3 \\
7 & 10 & 23.0 & 19.3 \\
8 & 10 & 16.4 & 30.0 \\
3 & 10 & 14.5 & 12.1 \\
\bottomrule
\end{tabular}
\end{table}

\subsection{Feature Policy}
\begin{table}[h]
\centering
\caption{Feature availability policy.}
\begin{tabular}{p{0.22\linewidth}p{0.68\linewidth}}
\toprule
Group & Policy \\
\midrule
Calendar & Known for the full forecast horizon. Used directly. \\
Historical targets & Used through lag, rolling, regime, and seasonal priors available at cutoff. \\
Promotions & Future realizations are not assumed. Only recurring historical priors are used. \\
Operations & Orders, traffic, returns, inventory, and reviews after 2022 are not used as realized inputs. \\
Mix features & Product, customer, source, and region mix are used only through train-derived priors. \\
\bottomrule
\end{tabular}
\end{table}

\subsection{Reproducibility Notes}
Primary scripts in the clean lineage include:
\begin{itemize}
  \item \texttt{run\_clean\_v10\_h1\_regime\_shape.py}
  \item \texttt{run\_clean\_v14\_multimetric\_frontier.py}
  \item \texttt{run\_clean\_v16\_foldlearned\_daily\_allocator.py}
  \item \texttt{run\_clean\_v17\_period\_ratio\_head.py}
  \item \texttt{run\_clean\_v19\_multimetric\_frontier.py}
\end{itemize}
Recommended checks before packaging are:
\begin{verbatim}
python -m py_compile run_clean_v19_multimetric_frontier.py
python run_clean_v19_multimetric_frontier.py
python run_multimetric_publiclike_research.py
\end{verbatim}

\subsection{External Analyst PDF}
The file \texttt{docs/datathon\_2026\_analyst.pdf} should be converted with Microsoft MarkItDown before final polishing:
\begin{verbatim}
markitdown "docs\datathon_2026_analyst.pdf" ^
  -o "docs\datathon_2026_analyst.md"
\end{verbatim}
Any claim from that document should be included only if it is verified against the provided datasets or existing clean logs.


\end{document}
